
\section{Intermediate Code Representation} 
\label{Sect_intermediate}

We introduce a format for representation of intermediate code.
For our type system, intermediate code cannot be generated in the
standard way.
Standard compilers create intermediate code immediately
after type checking. In the intermediate representation, all operations
have become unambiguous. This is possible because every variable
receives its type at the moment of initialization. 
Both LLVM (\cite{LLVM2004}) and GIMPLE (the intermediate representation
used by GCC, \cite{GIMPLE}) rely on unambiguous types for every value.
Unfortunately, this is not possible for our type system, 
because subtypes in general cannot be read off from declarations.
In the presence of loops, subtypes have to be obtained by means
of a fixed point computation. 

When an operation is ambiguous, all candidate overloads have to 
be inserted into the intermediate representation, after which 
those overloads whose preconditions are not met, need to be removed during
fixed point computation. 

\begin{definition}
   \label{Def_flatterm}
   We define \emph{flatterms}:
   \begin{itemize}
   \item
      If $ v $ is a local variable, then $ v $ is a flatterm.
   \item
      If $ c $ is a constant of a primitive type,
      then $ c $ is a flatterm.
   \item
      If $ {\cal A} $ is a DMA, and $ F $ is a state
      condition of $ {\cal A}, $ then $ {\bf some}({\cal A}, F ) $ is
      a flatterm.
   \item
      If $ v_1, \ldots, v_n $, are variables, $ f $ is an identifier,
      then $ f(v_1, \ldots, v_n) $ is a flatterm.
   \item
      If $ v $ is a variable, then
      $ {\bf len}(v), \ {\bf not}(v), \ {\bf neg}(v) $ 
      are flatterms.
   \item
      If $ v_1 $ and $ v_2 $ are variables,
      then 
      $ {\bf eq}(v_1,v_2), \ {\bf ne}(v_1,v_2), \
         {\bf lt}(v_1,v_2), \ {\bf gt}(v_1,v_2), $ \\
      $ {\bf le}(v_1,v_2), \ {\bf ge}(v_1,v_2), \
        {\bf add}(v_1,v_2), \ {\bf sub}(v_1,v_2)  $
      are flatterms.
  \item
      If $ v $ is a variable,
      $ {\cal A} $ is a DMA, and $ F $ is a state condition, then
      $ {\bf check}( v, {\cal A}, F ) $ is a flatterm.
   \item
          if $ v $ and $ w $ are variables,
          then $ {\bf ind}(v,w) $ is a flatterm.
   \item
      If $ v $ is a variable, and $ f $ is an identifier,
      then $ {\bf fld}_f(v) $ is a flatterm.
   \item
      If $ v_1, \ldots, v_n $ are variables with $ n \geq 0 $,
          then $ ( v_1, \ldots, v_n ) $ and
                   $ [ v_1, \ldots, v_n ] $ are flatterms.
   \end{itemize}
\end{definition}

\noindent
The meanings are as follows: If $ v $ is an array, then $ {\bf len}(v) $
equals its size. The intuitive meaning of $ {\bf ind}(v,w) $ is $ v[w]. $
The intuitive meaning of $ {\bf fld}_f(v) $ is $ v.f. $
Flatterm $ {\bf some}( {\cal A}, F ) $ non-deterministically creates a
data term $ t, $ s.t. $ F( \, Q_{\cal A}(t) \, ), $ while
$ {\bf check}( t, {\cal A}, F ) $ returns a boolean representing
truth of $ F( \, Q_{\cal A}(t) \, ). $
\begin{definition}
   \label{Def_program_point}
   We assume a countable set $ \Pi $ of program points.
\end{definition}
It is not very important how program points are represented.
Since statements defined below are trees, one can represent them as
paths.
\begin{definition}
   \label{Def_statement}
   We recursively define \emph{statements}:
   For a given statement, we assume that every substatement
   has a unique program point.
   A statement has one of the following forms:
   \begin{itemize}
   \item
      A scalar assignment of $ v \leftarrow t, $ in which $ v $ is a variable
      and $ t $ is a flatterm.
   \item
      A array assignment of $ v[w] := t $, in which
      $ v $ and $ w $ are variables,
      and $ t $ is a flatterm.
      Note that array assignment is a single operation
      with three arguments.
   \item
      If $ s_1, \ldots, s_n $ are statements, then
      $ \statblock(s_1, \ldots, s_n), \
        \statrepeat(s_1, \ldots, s_n), $ and 
        $ \statbranch(s_1, \ldots, s_n) $ are also statements.
   \item
      If $ v $ is a variable and
      $ s $ is a statement, then
      $ \statistrue \, v $ and $ \statisfalse \, v $ are also statements.
   \item
      $ \statnop $ is a statement.
  \item
      If $ v $ is a variable, then $ {\bf return} \, v  $ is a statement.
   \item
      $ v. {\rm push\_back}( w ) $
      with $ v,w $ variables is a statement.
   \item
      $ v. {\rm pop\_back}( ) $ is a statement.
   \item
      A field assignment $ w_1. f := w_2, $
      where $ f $ is a scalar field name,
      and $ w_1,w_2 $ are variables.
   \item
      If $ \overline{v} $ is a sequence of variables,
      and for each $ i $ with $ 1 \leq i \leq n, $ \ \
      $ {\cal A}_i $ is a sequence of automata,
      $ F_i $ is a sequence of state conditions,
      s.t. $ \| v \| = \| {\cal A}_i \| = \| F_i \|, $
      each $ s_i $ is a statement, then \\
      $ \statresolve( \overline{v}, ( {\cal A}_1, F_1 ) / s_1, \ldots,
                                    ( {\cal A}_n, F_n ) / s_n ) \, ) $
      is also a statement.
   \item
      If $ v $ is a variable, then $ \staterase \, v $ is a statement.
   \end{itemize}
\end{definition}

\noindent
Both $ \statrepeat $ and $ \statbranch $ are non-deterministic.
The $ \statrepeat $ statement non-deterministically skips
$ s_1, \ldots, s_n $ or starts at $ s_1. $
Similarly, $ \statbranch $ non-deterministically executes one of its
$ s_i. $ We chose for non-determinism because analyzing
non-deterministic code is not harder than deterministic code,
due to the fact that overapproximation will introduce non-determinism anyway.
In order to model deterministic code, one has to put
a $ \statistrue $ or $ \statisfalse $ at the beginning of every branch,
at the beginning of very \statrepeat, and immediately after every
\statrepeat.
Both statements quietly fail when the condition is not met.

In the $ \statresolve $ statement, the variables $ \overline{v} $ are used
for type checking. Overload $ ( {\cal A}_i, F_i ) / s_i $ is viable,
if for each variable in $ \overline{v}, $ the value, when
run with the corresponding automaton in $ {\cal A}_i, $ reaches a state
that satisfies the state corresponding condition.
\noindent
The difference between {\bf abort} and {\bf error} is
that $ {\bf error} $ is checked during type checking,
while $ {\bf abort} $ is checked at run time.

\begin{definition}
   \label{Def_state}
   A \emph{state} $ \sigma $ is a finite mapping from variables to
   data terms.
\end{definition}

We now define a function $ {\bf eval}_{\sigma}(t) $
that evaluates a flatterm $ t $ in state $ \sigma. $
Because the semantics of flatterms is non-deterministic,
$ {\bf eval}_{\sigma}(t) $ defines a set of possible values.

\begin{definition}
   \label{Def_semantics}
   For a given statement $ s, $ we assume two functions
   $ \Phi_{\bf in} $ and $ \Phi_{\bf out} $ that map
   program points to set of states.
   Both functions are defined for the program points occurring
   in $ s. $ They are defined together as the least fixed point
   satisfying the following conditions:
\begin{itemize}
\item
   If $ \pi $ contains a $ \statnop, $ then
   $ \Phi_{\bf in}( \pi ) \subseteq \Phi_{\bf out}( \pi ). $
\item
   If $ \pi $ contains $ \statblock( s_1, \ldots, s_n), $
   then let $ \pi_1, \ldots, \pi_n $ be the positions of
   $ s_1, \ldots, s_n. $
   We have $ \Phi_{\bf in}( \pi ) \subseteq \Phi_{\bf in}( \pi_1). $
   For every $ i $ with $ 1 \leq i \leq n-1, $ we have
   $ \Phi_{\bf out}( \pi_i ) \subseteq \Phi_{\bf in}( \pi_{i+1} ), $ and
   $ \Phi_{\bf out}(\pi_n) \subseteq \Phi_{\bf out}(\pi). $
\item
   If $ \pi $ contains $ \statrepeat( s_1, \ldots, s_n ), $ then
   let $ \pi_1, \ldots, \pi_n $ be the positions of
   $ s_1, \ldots, s_n. $
   We have $ \Phi_{\bf in}( \pi ) \subseteq \Phi_{\bf in}( \pi_1). $
   For every $ i $ with $ 1 \leq i \leq n-1, $ we have
   $ \Phi_{\bf out}( \pi_i ) \subseteq \Phi_{\bf in}( \pi_{i+1} ). $
   In addition we have $ \Phi_{\bf in}(\pi) \subseteq \Phi_{\bf out}(\pi), $ and
   $ \Phi_{\bf out}( \pi_n ) \subseteq \Phi_{\bf in}(\pi). $
\item
   If $ \pi $ contains $ \statbranch( s_1, \ldots, s_n ), $
   then let $ \pi_1, \ldots, \pi_n $ be the positions of
   $ s_1, \ldots, s_n. $
   For every $ i, $ we have
   $ \Phi_{\bf in}( \pi ) \subseteq \Phi_{\bf in}( \pi_i ), $ and
   $ \Phi_{\bf out}( \pi_i ) \subseteq \Phi_{\bf out}(\pi). $
\item
   If $ \pi $ contains $ \statistrue \, v, $ then
   $ \Phi_{\bf in}( \pi ) \cap
     \{ \sigma \mid \sigma(v) = {\bf t} \} \subseteq
     \Phi_{\bf out}( \pi). $
\item
   $ \statisfalse $ is analogous to $ \statistrue. $
\item
   If $ \pi $ contains a statement of form $ v \leftarrow t $ with
   $ t $ a flatterm, then
   \[ \{ \sigma \cup \{ v := d \} \, \} \mid
           \sigma \in \Phi_{\bf in}(\pi) \mbox{ and }
           d \in {\bf eval}_{\sigma}(t) \, \} \subseteq
           \Phi_{\bf out}( \pi ). \]
   Function $ {\bf eval}_{\sigma}(t) $ is a function
   that (nondeterministically) evaluates flatterms. 
\item
   If $ \pi $ contains $ \statresolve( c_1/s_1, \ldots, c_n/s_n ), $
   then let $ \pi_1, \ldots, \pi_n $ be the positions of
   $ s_1, \ldots, s_n. $ Similary, let $ \rho_1, \ldots, \rho_n $
   be the positions of $ c_1, \ldots, c_n. $
   We have $ \Phi_{\bf in}(\pi) \subseteq \Phi_{\bf in}( \pi_i ) $ and
   $ \Phi_{\bf in}( \pi ) \subseteq \Phi( \rho_i ). $

   We also have
   \[ \bigcap \Phi_{\bf out}( \pi_i ) \subseteq \Phi_{\bf out}( \pi ). \]
\end{itemize}
\end{definition}
The last condition is a bit tricky.
It is possible that some of the $ s_i $ use temporary
variables that other $ s_i $ do not use.
In that case, the resulting state will always be unique,
and not in the intersection. This problem can be solved by
erasing temporary variables. In fact, this is the reason
that $ \staterase $ was created.

We do not store any scope
information for local variables in statements.
This implies that local variables
in principle remain visible, once initialized.
This may lead to
space inefficiency when statements are naively translated into $ C $ or
$ C^{++}, $ which allow more restricted scopes. This
problem can be solved by
doing a stack based liveness analysis before translation.
\begin{definition}
   \label{Def_live}
   We define $ {\bf live}_{\bf in}(\pi) $ and
   $ {\bf live}_{\bf out}(\pi) $ as the least
   fixed points of the propagation rules below. The intuitive
   meaning is: $ {\bf live}(\pi)_{\bf in/out} $ are
   the variables whose value may be later used in a computation
   that passes through $ \pi. $
   \begin{itemize}
   \item
      If $ \pi $ contains a statement of form
      $ v \leftarrow t, $ then let $ V $ be the
      set of variables occurring in flatterm $ t. $
      \[ \Phi_{\bf in}(\pi) =
         ( \Phi_{\bf out}(\pi) \backslash \{ v \} ) \cup V. \]
   \item
      If $ \pi $ contains a statement of form
      $ \statisfalse \, v, \ \statistrue \, v, \ \statreturn \, v $
      then
      \[ {\bf live}_{\bf in}( \pi ) =
                {\bf live}_{\bf out}( \pi ) \cup \{ v \}. \]

   \item
      If $ \pi $ contains a statement of form
      $ v. {\rm push\_back}( w ), $ then
      \[ {\bf live}_{\bf in}(\pi) = {\bf live}_{\bf out}(\pi) \cup
              \{ w, v \}. \]
   \item
      If $ \pi $ contains $ \statnop, $ then
      $ {\bf live}_{\bf out}( \pi ) \subseteq {\bf live}_{\bf in} ( \pi ). $
   \item
      If $ \pi $ contains $ \statbranch( s_1, \ldots, s_n ), $
      then let $ \pi_1, \ldots, \pi_n $ be the positions of
      $ s_1, \ldots, s_n. $
      For every $ i, $ we have
      $ \Phi_{\bf in}(\pi_i) \subseteq \Phi_{\bf in}( \pi ), $ and
      $ \Phi_{\bf out}( \pi ) \subseteq \Phi_{\bf out}(\pi_i). $

   \end{itemize}
  A variable $ v $ is \emph{live} at some program point $ \pi $
   if $ v \in {\bf live}_{\bf in}(\pi). $
\end{definition}
Instead of defining $ {\bf live}(\pi) $ as a set of variables, one can also
define $ {\bf use}(\pi) $ as a set of pairs $ ( v, \pi'), $
indicating that $ \pi' $ is a point where the value of $ v $ is potentially used.
The definition is only slightly more complicated than Definition~\ref{Def_live} above.

